\section{Discussion}
\label{sec:discussion}

\subsection{Methodological Caveats}
\label{sec:discussion-caveats}

\begin{enumerate}
  \item \textbf{Pearson vs.\ rank correlations.} The pipeline currently
  uses Pearson correlation, which is sensitive to outliers and to
  non-Gaussian marginals.  Spearman rank correlation is offered in
  the code; results are qualitatively similar but the community
  partition differs in detail.

  \item \textbf{Choice of distance.} The Mantegna distance
  \eqref{eq:mantegna} is the de-facto standard in the financial-network
  literature but is not the only choice.  Alternative metrics such as
  the angular (cosine) distance used by the PMFG literature or the
  log-distance $d_{ij}=-\ln|\rho_{ij}|$ yield different topologies;
  we have not performed a head-to-head comparison and leave it for
  future work.

  \item \textbf{Louvain non-determinism.} The Louvain algorithm is
  stochastic~\cite{blondel2008fast}.  Our $k$-sweep reports
  $\bar Q \pm \sigma_Q$ across ten seeds and the operational code
  averages three deterministic seeds; we recommend that any downstream
  use of the partition for trading signal generation be re-seeded at
  every re-fit.

  \item \textbf{Modularity resolution limit.} Modularity maximisation
  has a well-known resolution limit~\cite{fortunato2007resolution}: it
  tends to merge small, closely-related communities.  The drop from
  $16$ to $6$ communities as $k$ grows from $2$ to $20$
  (Figure~\ref{fig:ksweep-modularity}B) is partly attributable to
  this, not to a real absence of structure.

  \item \textbf{Stationarity.} The Pearson correlation assumes a
  stationary covariance structure over the window.  The
  \emph{lookback} toggle ($1$/$3$/$5$/\emph{max} years) allows the
  user to probe the sensitivity of the partition to this assumption;
  the longer windows tend to be dominated by long-run co-movement, the
  shorter by recent regime dynamics.
\end{enumerate}

\subsection{Current Limitations}
\label{sec:discussion-limitations}

Beyond the methodological caveats of Section~\ref{sec:discussion-caveats},
we acknowledge four practical limitations of the present study that any
downstream user of the framework should bear in mind.
\begin{itemize}
  \item \textbf{Survivorship bias.}  The dataset comprises only stocks
  that \emph{survived} to the present day; bankrupt, suspended, and
  delisted companies are absent.  Consequently the portfolio
  backtest of Section~\ref{sec:results-portfolio} benefits slightly
  from hindsight, because the realised equity curve is constructed on
  a survivor-only universe rather than the universe of stocks that
  traded at each historical date.  A fully survivorship-bias-free
  treatment would require delisting-event data with the historical
  reason for exit, which is not uniformly available for NEPSE prior
  to 2010.

  \item \textbf{Illiquidity and zero-return bias.}  NEPSE is a
  frontier market in which many micro-cap stocks trade on a small
  fraction of all sessions.  Our missing-data policy of
  forward-filling the last observed price on non-trading days produces
  strings of exactly zero log returns.  Two highly illiquid tickers
  will exhibit near-zero recentred returns simultaneously and hence
  carry an artificially inflated Pearson correlation---a known
  artefact of zero-return filling in illiquid universes.  Practitioners
  should treat the correlation matrix as a \emph{lower bound} on the
  genuine covariance, and treat the resulting anomaly score $\eta$ as
  conservative.

  \item \textbf{Macro-stationarity.}  Compressing almost 23~years of
  data into a single \emph{max}-window correlation graph implicitly
  assumes that pairwise stock relationships are static across the
  sample.  In reality, NEPSE's regulatory regime has shifted through
  major episodes (NRB capital-hike policies, hydro-power IPP
  regime changes, the 2015 trade blockade, the 2020 pandemic shock),
  any of which can rearrange the correlation structure at the
  portfolio level.  We partially mitigate this in
  Section~\ref{fig:community-stability} via the rolling 1-year ARI
  / NMI stability test, which exposes non-stationarity explicitly,
  but the headline \emph{max}-window graph still compresses decades of
  macroeconomic regime-shifts into a single partition.

  \item \textbf{Market impact and slippage.}  The backtest of
  Section~\ref{sec:results-portfolio} assumes a flat 50~bps per-side
  transaction cost.  In reality, rebalancing a portfolio into thinly
  traded NEPSE names would interact with the order book's depth and
  produce significant price impact, especially in low-capitalisation
  segments.  Reported excess returns over the equal-weight benchmark
  are therefore an \emph{upper bound} on what a real trader could
  realise; practical execution would erode them by an unknown amount
  that depends on portfolio size and order-routing strategy.
\end{itemize}

\subsection{Practical Implications}
\label{sec:discussion-implications}

The community partition produced by NepGraph is a data-driven
complement to NEPSE's official sector taxonomy.  We see three concrete
use cases:
\begin{itemize}
  \item \textbf{Risk management.} A portfolio concentrated in a single
  detected community is exposed to a single common factor, even when
  the constituents span multiple official sectors.  The diversification
  score $\Delta$ is a one-number summary of this exposure.

  \item \textbf{Surveillance.} The anomaly score $\eta$ flags
  cross-sector MST links that may indicate coordinated price
  movement, contagion channels, or regulatory arbitrage.  We
  deliberately do \emph{not} claim that all flagged stocks are involved
  in market manipulation; we provide a list of candidates that
  warrants further qualitative investigation.

  \item \textbf{Index construction.} The detected communities can be
  used to construct \emph{de-correlated} equal-weight baskets, which
  the dashboard exposes via the backtest tab.
\end{itemize}

\subsection{Reproducibility}
\label{sec:discussion-reproducibility}

All code is open-source (MIT).  The $k$-sweep requires only the
bundled \texttt{nepse\_prices.csv} (~12\,MB) and runs in under five
minutes on a laptop.  The dashboard can be started with
\texttt{streamlit\ run\ dashboard.py} and updates the underlying MST,
Louvain partition and anomaly list interactively as the user toggles
the lookback or the top-$k$ slider.

\newpage
